\printedpage{24}
% head: A. Grothendieck : Dérivateurs, Ch. I, Édition des Universités de Montpellier II et Paris VII

\noindent\textbf{Der 2} (Second localization axiom) : For every $I \in \Dia$, the canonical functor
\[ \D(I) \longrightarrow \underline{\mathsf{Hom}}(I^{\circ}, \mathcal{A}_{\D}) \quad , \]
$\mathcal{A}_{\D} = \D(e)$ fundamental category for the derivator $\D$, is conservative.\\
Stable under induction, and under duality.

\noindent\msp{page 44}\texttt{[page 44]}

\noindent A prederivator is called a \textit{left derivator} if it satisfies \textbf{Der 1} and \textbf{Der 2}, and in addition the following axioms (of which here is the first).

\noindent\textbf{Der 3} (g) (Existence of direct images) : For every $f : I \to J$ in $\Dia$, the functor $f^{*} : \D(J) \to \D(I)$ admits a \textit{right adjoint} $f_{*}$\,.\\
Stable under induction.

We call $f_{*}$ the direct image functor (cohomological) for $f$. The object $f_{*}(F) \in \D(J)$ (for $F \in \D(I)$) is also called the object of ``\textit{relative cohomology}'' of $I$ over $J$ (via $f$), with coefficients in $F$. We have canonical isomorphisms
\[ (gf)_{*} \simeq g_{*} f_{*} \quad . \]

\noindent\textbf{Corollary.} \textit{The categories $\D(I)$ have finite products, the functors $f^{*}$ and $f_{*}$ commute with them.}

Uses \textbf{Der 1} for $I = \emptyset$ and $I = e \amalg e$, and \textbf{Der 3} for $I \xrightarrow{f} e$.

We have a dual axiom :

\noindent\textbf{Der 3} (d) : existence of left adjoints $f_{!}$ to the $f^{*}$, the functor $f_{!}$ is called ``homological direct image'' via $f$, and we have a notion of \textit{relative homology} of $I$ over $J$ (for $f$), with coefficients in $F$.

In what follows, the indication (g) or (d) next to the designation \textbf{Der} $n$ of an axiom, % [sic] coté -> côté, missing circumflex
means that it applies to left derivators, respectively to right derivators. Likewise for statements, definitions \textit{etc.}

\noindent\msp{page 45}\texttt{[page 45]}

\noindent\textbf{Der 4}$_{0}$ (g)\footnote[13]{Unnecessary as an axiom, will become a theorem (follows from the other axioms).} : Let $U \overset{j}{\hookrightarrow} X \overset{i}{\hookleftarrow} Y$ be in $\Dia$, open and closed immersions complementary to one another. Then the adjunction morphism in $\D(Y)$
\[ i^{*} i_{*}(F_{Y}) \longrightarrow F_{Y} \]
is an isomorphism, hence
\[ i_{*} : \D(Y) \longrightarrow \D(X) \]
is fully faithful, and $i^{*} : \D(X) \to \D(Y)$ is a localization functor. The essential image of $i_{*}$ consists of the objects $F$ of $\D(X)$ such that $j^{*}(F)$ is the final object of $\D(U)$.

% foot: Version 8 juillet 2001

\printedpage{25}
% head: A. Grothendieck : Dérivateurs, Ch. I, Édition des Universités de Montpellier II et Paris VII

\noindent\textbf{Der 4}$_{0}$ (d) : With the preceding notation for $i$, $j$, the adjunction morphism in $\D(U)$
\[ F_{U} \longrightarrow j^{*} j_{!} F_{U} \]
is an isomorphism, - $j_{!}$ is fully faithful, $j^{*}$ a localization functor. The essential image of $j_{!}$ is the full subcategory of $\D(X)$ consisting of the $F$ such that $i^{*}(F)$ is an initial object of $\D(Y)$.

There are the two squares of page 8 and 9, but as for interpreting the cartesian or cocartesian condition it is premature here.

\noindent\msp{page 46}\texttt{[page 46]}

\noindent Base change axioms.

\noindent\textbf{Der 4} (g) (Fibers of a direct image) : Let $f : X \to Y$ be in $\Dia$, $y \in Y$, let us consider the known diagram
\[
\begin{tikzcd}[row sep=large, column sep=large]
X/y \arrow[r, "q"] \arrow[d, "f'"'] & X \arrow[d, "f"] \\
e \arrow[r, "p = \alpha_{y}"'] & Y
\end{tikzcd}
\]
and the corresponding base change morphism
\[ f_{*}(F)_{y} \overset{\text{def}}{=} p^{*} f_{*}(F) \longrightarrow f'_{*}(q^{*}(F))\bigl(\overset{\text{def}}{=} \Gamma_{\D}(X/y, q^{*}(F))\bigr) \quad . \]

\noindent This is always an isomorphism.

\noindent\msp{page 47}\texttt{[page 47]}

\noindent\textbf{Der 5} (g) (Cofinality axiom for cohomology)\footnote[14]{It is a consequence of weaker statements.} : Let $\alpha : X' \to X$ in $\Dia$ be ``cohomologically cofinal'' \textit{i.e.} the $x\backslash X'$ are contractible ($x \in X$) \textit{i.e.} $\alpha^{\circ} : {X'}^{\circ} \to X^{\circ}$ is a Hot-equivalence ``locally over $X^{\circ}$\,''. Then $\forall F \in \D(X)$, the homomorphism
\[ \Gamma_{\D}(X, F) \longrightarrow \Gamma_{\D}(X', \alpha^{*} F) \]
is an isomorphism (in $\mathcal{A} = \D(e)$).

\noindent N.B. If $f : X \to e$, we set $\Gamma_{\D}(X, F) = f_{*}(F)$ (``absolute'' cohomology of $X$, with coefficients in $F$).

\noindent Same thing\footnote[15]{False!} if $\alpha$ is ``homologically cofinal'', i.e. the $X'/x$ are contractible \texttt{[sentence crossed out]}.

% foot: Version 8 juillet 2001

\printedpage{26}
% head: A. Grothendieck : Dérivateurs, Ch. I, Édition des Universités de Montpellier II et Paris VII

\noindent\textbf{Remark.} We should be able to strengthen this axiom considerably, by requiring that
\[ \alpha_{!}(\varepsilon_{X'}) \xrightarrow{\sim} \varepsilon_{X} \qquad\qquad \text{in} \quad \HOT(X) \quad , \]
where $\varepsilon_{X'}$, $\varepsilon_{X}$ \texttt{[are]} the final objects of $\HOT(X')$, $\HOT(X)$, and where $\alpha_{!}(\varepsilon_{X'}) \to \varepsilon_{X}$ is the adjunction homomorphism, taking into account that $\varepsilon_{X'} \simeq \alpha^{*}(\varepsilon_{X})$. We say that $X'$ is \textit{homotopically trivial over $X$}. (And if the dual condition is satisfied in $\HOT(X^{\circ})$, we say that $X'$ is cohomotopically trivial over $X$.) This condition is probably \underline{not} self-dual. It is satisfied both if $\alpha$ is cohomologically cofinal, and if it is homologically cofinal (costrict or strict Hot-equivalence).

\noindent N.B. For $f : X \to e$, saying that $f$ is cohomologically cofinal, homologically cofinal, or that $X$ is contractible, is equivalent.

\noindent\textbf{Corollary 1} (g) (Asphericity theorem for constant coefficients in cohomology). \textit{Let $f : X \to e$ be in} $\Dia$\textit{, $X$ contractible, $M \in \D(e)$. Then the adjunction homomorphism}
\[ M \xrightarrow{\sim} f_{*} f^{*}(M) \quad \bigl(\overset{\text{def}}{=} \Gamma(X, f^{*}(M))\bigr) \]
\textit{is an isomorphism.}

\noindent\textbf{Corollary 2} (g)\footnote[16]{\underline{False}, except for constant coefficients.} \textit{Let $f : X \to Y$ be} \texttt{[a]} \textit{strict Hot-equivalence (i.e. $\forall y \in Y$, $X/y$ contractible). Then for $F \in \Ob\,\D(Y)$}
\[ F \overset{\text{adj}}{\xrightarrow{\sim}} f_{*} f^{*} F \]
\texttt{[is an]} \textit{isomorphism.} \texttt{[This corollary is crossed out, but not its proof.]}

\noindent\msp{page 48}\texttt{[page 48]}

\noindent We apply \textbf{Der 2}, \textbf{Der 4} (g), to reduce to looking, for $y_{0} \in Y$, at the homomorphism $\alpha$ of the following (commutative!) diagram :
\[
\begin{tikzcd}[row sep=large, column sep=huge]
\Gamma(\{y_{0}\}, F_{y_{0}}) = F_{y_{0}} \arrow[r, "\alpha"] & \Gamma_{\D}(X/y_{0}, f^{*}(F)) \\
\Gamma(Y/y_{0}, F) \arrow[u, "\varphi\,\wr"] \arrow[ur, "\alpha'"'] &
\end{tikzcd}
\quad ,
\]
where $\alpha'$ is deduced from the functor induced by $f$
\[ X/y_{0} \longrightarrow Y/y_{0} \]
and where $\varphi$ is deduced from the functor
\[ \{y_{0}\} \longrightarrow Y/y_{0} \qquad \text{(final object)}\,. \]

% foot: Version 8 juillet 2001

\printedpage{27}
% head: A. Grothendieck : Dérivateurs, Ch. I, Édition des Universités de Montpellier II et Paris VII

\noindent Thus $\varphi$ is an isomorphism by \textbf{Der 5} (g)\footnote[17]{\textit{cf.} corollary 3, page 50, which we should move up here.} (N.B. If $z \in Z$, whence $\{z\} \to Z$, this functor is cohomologically cofinal if and only if $z$ is a final object, homologically final \texttt{[rather cofinal]} if and only if $z$ is an initial object). Thus the conclusion of the corollary is equivalent to saying that $\alpha'$ is an isomorphism. But $X/y_{0} \to Y/y_{0}$ is also homologically cofinal, for if $u : y_{1} \to y_{0}$ is in $Y/y_{0}$, we have a canonical isomorphism
\[ (X/y_{0})/(y_{1}, u) \simeq X/y_{1} \]
which is contractible. So by the second part of \textbf{Der 5} (g),\footnote[18]{(which is \underline{false})} we win.

\noindent\msp{page 49}\texttt{[page 49]}

\noindent\textbf{Corollary} 4 (g) (base change criterion for cohomological direct images). \textit{Let}
\[
\begin{tikzcd}[row sep=huge, column sep=huge]
X' \arrow[r, "q"] \arrow[d, "f'"'] & X \arrow[d, "f"] \arrow[dl, dashed, "\gamma"', shorten <=4ex, shorten >=1ex] \\
Y' \arrow[r, "p"'] & Y
\end{tikzcd}
\]
\textit{be in} $\Dia$ (\textit{cf.} page 30), \textit{whence the base change homomorphism}
\[ p^{*} f_{*}(F) \longrightarrow f'_{*} q^{*}(F) \qquad (F \in \Ob\,\D(X)) \quad . \]
\textit{For this to be an isomorphism for all $F$, it is enough that for every $y' \in \Ob\,Y'$, the functor induced by} (p. 9)
\[ X'/y' \longrightarrow X/p(y') \]
\textit{be a strict Hot-equivalence (or costrict);\footnote[19]{\underline{False} for ``strict'', works \underline{only} for ``\underline{\underline{co}}strict''.} this also means (in the strict case) that for}
\[ y'_{0} \in \Ob\,Y' \,,\ \ x_{0} \in \Ob\,X \,,\ \ u_{0} : f(x_{0}) \longrightarrow p(y'_{0}) \]
\textit{the category $K(y'_{0}, x_{0}, u_{0})$ of objects of the form}
\[ (x', v', u) \mid x' \in \Ob\,X',\ \ v' : f'(x') \to y'_{0},\ \ u : q(x') \overset{\dashleftarrow}{\longrightarrow} x_{0} \]
\textit{making the square in $Y'$ commutative}
\[
\begin{tikzcd}[row sep=huge, column sep=huge]
f(q(x')) \arrow[r, "\gamma(x')"] \arrow[d, shift right=0.5ex, "f(u)"'] & p(f'(x')) \arrow[d, "p(v')"] \\
f(x_{0}) \arrow[r, "u_{0}"'] \arrow[u, dashed, shift right=0.5ex, "\leftarrow\ \text{\scriptsize\shortstack[l]{case of\\costrict\\Hot-equivalence}}"'] & p(y'_{0})
\end{tikzcd}
\tag*{$(*)$}
\]
\textit{is contractible.}

% foot: Version 8 juillet 2001

\printedpage{28}
% head: A. Grothendieck : Dérivateurs, Ch. I, Édition des Universités de Montpellier II et Paris VII

\noindent N.B. In the case of the \underline{co}strict Hot-equivalence, it is the same description of
\[ K = (x_{0}, u_{0})\backslash(X'/y'_{0}) \quad , \]
except that the arrow $u$ is in the opposite direction $x_{0} \to q(x')$ and likewise in the commutative diagram $(*)$.\\
We apply \textbf{Der 2}, \textbf{Der 4} (g) and corollary 1 \texttt{[rather} \textbf{Der 5} (g)\texttt{]}.

\noindent\msp{page 50}\texttt{[page 50]}

\noindent\textbf{Corollary 5} (g). \textit{Let $f : X \to Y$ be fully faithful in} $\Dia$\textit{, and $F \in \D(X)$. Then the adjunction homomorphism}
\[ f^{*} f_{*}(F) \longrightarrow F \]
\textit{is an isomorphism.}

Indeed, it is enough to check it fiber by fiber (\textbf{Der 2}). Let $x \in X$, $y = f(x) \in Y$, % [sic] efffet -> effet, tripled f
then
\[ f^{*} f_{*}(F)_{x} \simeq f_{*}(F)_{y} \simeq \Gamma_{\D}(X/y, F) \]
now because of full faithfulness $X/y \simeq X/x$.

\noindent\textbf{Corollary 3} (move up to page 48).\footnote[20]{Will follow from the stronger \textbf{Der 6} (g) below.} \textit{Let $Z \in$} $\Dia$\textit{, suppose that $Z$ has a final object $e_{Z}$. Let $p : e \to Z$ be the corresponding functor. Then the canonical homomorphism (for $F$ in} $\D(Z)$\,)
\[
\begin{tikzcd}[row sep=huge]
Z \arrow[d, shift left=0.6ex, "f"] \\
e \arrow[u, shift left=0.6ex, "p"]
\end{tikzcd}
\qquad\qquad\qquad\qquad
f_{*}(F) \longrightarrow p^{*}(F)
\]
\textit{is an isomorphism.}\footnote[21]{Axiom of the index category with a final object (arbitrary coefficient).}

This homomorphism is \texttt{[a]} special case of this one, for
\[
\begin{tikzcd}[row sep=large, column sep=large]
Z' \arrow[rr, "p"] \arrow[dr, "f'"'] & & Z \arrow[dl, "f"] \\
& S &
\end{tikzcd}
\]
(taking $Z' = S = e$, $f' = \mathsf{id}_{e}$),
\[ f_{*}(F) \longrightarrow f'_{*}(p^{*}(F)) \quad , \qquad (\, F \in \Ob\,\D(Z)\,) \quad , \]
and it is a special case\footnote[22]{already considered in cor. 2, p. 47} of the morphisms of

% foot: Version 8 juillet 2001